\documentclass[11pt,a4paper]{article}

\usepackage{fontspec}
\setmainfont{Latin Modern Roman}
\newfontfamily{\hindifont}{Nirmala UI}[Script=Devanagari]
\newcommand{\hindi}[1]{{\hindifont #1}}

\usepackage[top=1.8cm,bottom=2cm,left=2cm,right=2cm]{geometry}
\usepackage{xcolor}
\usepackage{tcolorbox}
\usepackage{enumitem}
\usepackage{booktabs}
\usepackage{tabularx}
\usepackage{fancyhdr}
\usepackage{wasysym}
\usepackage{multicol}
\usepackage{wasysym}

\definecolor{govblue}{HTML}{003566}
\definecolor{govgold}{HTML}{c9a84c}
\definecolor{lightgray}{HTML}{f5f5f5}
\definecolor{streama}{HTML}{e8f4fd}
\definecolor{streamb}{HTML}{eafaf1}

\pagestyle{fancy}
\fancyhf{}
\fancyhead[L]{\small\color{govblue}\textbf{AIGGPA Field Survey -- IT Utilisation Assessment}}
\fancyhead[R]{\small\color{govblue}April 2026}
\fancyfoot[C]{\small\thepage}
\renewcommand{\headrulewidth}{0.4pt}
\setlength{\headheight}{14pt}

\newcommand{\q}[2]{%
  \vspace{4pt}
  \noindent\textbf{#1.}~#2\\[2pt]
}

\newcommand{\ans}{%
  \noindent\rule{\textwidth}{0.3pt}\\[3pt]
}

\newcommand{\choices}[1]{%
  \begin{itemize}[label=$\square$, leftmargin=1.4em, itemsep=1pt, topsep=2pt]
    #1
  \end{itemize}
}

\begin{document}

%===== COVER BOX =====
\begin{tcolorbox}[colback=govblue,colframe=govblue,arc=2mm]
  \color{white}
  \centering
  {\large\bfseries IT Utilisation Assessment -- Field Survey Instrument}\\[0.3em]
  {\normalsize\bfseries \hindi{आईटी उपयोग आकलन -- क्षेत्र सर्वेक्षण प्रपत्र}}\\[0.5em]
  {\small Atal Bihari Vajpayee Institute of Good Governance \& Policy Analysis (AIGGPA), Bhopal}\\
  {\small \textit{Strictly for academic research. Responses are confidential and anonymous.}\\
  \hindi{यह प्रपत्र शैक्षणिक अनुसंधान हेतु है। उत्तर पूर्णतः गोपनीय हैं।}}
\end{tcolorbox}

\vspace{0.4cm}

%===== RESPONDENT DETAILS =====
\begin{tcolorbox}[colback=lightgray,colframe=govblue!50,arc=2mm,boxrule=0.5pt,title=\textbf{Respondent Details \hindi{(प्रतिभागी विवरण)}} -- \textit{not shared publicly},fonttitle=\small]
\small
\noindent
\begin{tabularx}{\textwidth}{@{}XX@{}}
  \textbf{Department \hindi{(विभाग)}:} \rule{5cm}{0.3pt} &
  \textbf{Office Level \hindi{(कार्यालय स्तर)}:} \\[2pt]
  & \choices{
    \item State / \hindi{राज्य}
    \item District / \hindi{जिला}
    \item Block / \hindi{ब्लॉक}
  }\\[4pt]
  \textbf{Designation \hindi{(पद)}:} \rule{5cm}{0.3pt} &
  \textbf{Grade \hindi{(श्रेणी)}:}
  \choices{
    \item Class I (Group A)
    \item Class II (Group B)
    \item Class III (Group C)
  }\\[4pt]
  \textbf{Years in service \hindi{(सेवा वर्ष)}:} \rule{3.5cm}{0.3pt} &
  \textbf{Date of survey:} \rule{3.5cm}{0.3pt}
\end{tabularx}
\end{tcolorbox}

\vspace{0.5cm}

%=========================================================
\begin{tcolorbox}[colback=streama,colframe=govblue,arc=2mm,boxrule=0.6pt,
  title={\textbf{STREAM A -- For Department Officials \& Employees} \quad \small\hindi{(विभागीय अधिकारी एवं कर्मचारी)}},fonttitle=\normalsize]

%--- PART 1: ROLES ---
\textbf{\color{govblue}Part 1 -- Your Role \& Daily Tasks \hindi{(आपकी भूमिका एवं दैनिक कार्य)}}
\vspace{4pt}

\q{1}{What are your 3 main responsibilities in this role? \hindi{(इस पद में आपकी 3 मुख्य जिम्मेदारियां क्या हैं?)}}
\ans\ans\ans

\q{2}{Which of your daily tasks do you think \textbf{could benefit from} an IT tool or computer system?\\
\hindi{आपके किन दैनिक कार्यों को आईटी उपकरण या कंप्यूटर से सहायता मिल सकती है?}}
\ans\ans

%--- PART 2: CURRENT IT TOOL USE ---
\vspace{4pt}
\textbf{\color{govblue}Part 2 -- IT Tools You Currently Use \hindi{(वर्तमान में उपयोग किए जाने वाले आईटी उपकरण)}}
\vspace{4pt}

\q{3}{Which of the following IT tools do you personally use for your work? (Tick all that apply)\\
\hindi{आप अपने कार्य के लिए निम्नलिखित में से कौन से आईटी उपकरण उपयोग करते हैं?}}
\begin{multicols}{2}
\choices{
  \item e-Office / e-File \hindi{(ई-कार्यालय)}
  \item HRMIS / Service Book \hindi{(सेवा रजिस्टर)}
  \item SPARROW (online APAR)
  \item iGOT Karmayogi
  \item Department Portal (specify): \rule{1.5cm}{0.3pt}
  \item MS Excel / Google Sheets
  \item MS Word / Google Docs
  \item WhatsApp (for official work)
  \item Email (official)
  \item Dashboard / MIS Report system
  \item None of the above
  \item Other: \rule{2cm}{0.3pt}
}
\end{multicols}

\q{4}{For writing official letters and communications -- how do you currently do it? \hindi{(आधिकारिक पत्र लिखने के लिए आप क्या तरीका उपयोग करते हैं?)}}
\choices{
  \item I type it myself on a computer / laptop
  \item A typist or data entry operator types it for me
  \item I handwrite it and it is typed by someone else
  \item I use a pre-prepared template on the computer
  \item Other: \rule{6cm}{0.3pt}
}

\q{5}{For which specific tasks do you use \textbf{dashboards, reports, or data systems} in your role?\\
\hindi{आप अपनी भूमिका में डैशबोर्ड, रिपोर्ट या डेटा प्रणाली किन कार्यों के लिए उपयोग करते हैं?}\\
If not applicable, write ``Not applicable''.}
\ans\ans

%--- PART 3: NON-USE & BARRIERS ---
\vspace{4pt}
\textbf{\color{govblue}Part 3 -- Why Tools Are Not Used \hindi{(उपकरण उपयोग न करने के कारण)}}
\vspace{4pt}

\q{6}{If there are tasks in your role where you \textbf{do not use} an IT tool but feel one would help -- why have you not used it? (Tick all that apply)\\
\hindi{यदि ऐसे कार्य हैं जिनमें आप आईटी उपकरण उपयोग नहीं करते, तो क्यों नहीं?}}
\choices{
  \item I was never trained on it \hindi{(प्रशिक्षण नहीं मिला)}
  \item No computer or device is available at my desk \hindi{(कंप्यूटर उपलब्ध नहीं)}
  \item No internet connection available \hindi{(इंटरनेट उपलब्ध नहीं)}
  \item I am not aware of which tool to use \hindi{(कौन सा उपकरण उपयोग करूं, पता नहीं)}
  \item The manual method is faster/easier for me \hindi{(मैनुअल तरीका आसान लगता है)}
  \item The system is often not working / slow \hindi{(सिस्टम अक्सर काम नहीं करता)}
  \item My seniors do not require the use of IT tools \hindi{(वरिष्ठ अधिकारियों की आवश्यकता नहीं)}
  \item Other: \rule{6cm}{0.3pt}
}

%--- PART 4: LEARNING ---
\vspace{4pt}
\textbf{\color{govblue}Part 4 -- How Did You Learn to Use IT Tools \hindi{(आईटी उपकरण सीखने का तरीका)}}
\vspace{4pt}

\q{7}{How did you learn to use the IT tools you currently use? (Tick all that apply)\\
\hindi{आपने जो आईटी उपकरण उपयोग करते हैं, वे आपने कैसे सीखे?}}
\choices{
  \item Formal government training / workshop \hindi{(सरकारी प्रशिक्षण कार्यशाला)}
  \item Trained by a colleague at the office \hindi{(कार्यालय सहकर्मी द्वारा)}
  \item Self-taught / YouTube / internet \hindi{(स्व-अध्ययन)}
  \item Learned in school / college before joining service \hindi{(सेवा पूर्व शिक्षा में)}
  \item I have not learned / do not use any IT tools \hindi{(मैं कोई आईटी उपकरण नहीं सीखा)}
}

\q{8}{Would you like to receive formal training on any IT tool? If yes, which one?\\
\hindi{क्या आप किसी आईटी उपकरण पर औपचारिक प्रशिक्षण प्राप्त करना चाहेंगे? यदि हां, तो कौन सा?}}
\choices{
  \item Yes -- \hindi{हां:} \rule{8cm}{0.3pt}
  \item No, I am comfortable with current tools \hindi{(नहीं, मैं वर्तमान उपकरणों से संतुष्ट हूं)}
}

\end{tcolorbox}

\vspace{0.5cm}

%=========================================================
\begin{tcolorbox}[colback=streamb,colframe=govblue!60,arc=2mm,boxrule=0.6pt,
  title={\textbf{STREAM B -- For IT Department / IT Cell Officials} \quad \small\hindi{(आईटी विभाग / प्रकोष्ठ के अधिकारी)}},fonttitle=\normalsize]

\textit{\small This stream is for officials working in or heading the IT unit/section of the department. Stream A questions may be skipped for these respondents.}
\vspace{6pt}

\q{B1}{What IT tools, platforms, or software does your IT department/cell currently operate and maintain?\\
\hindi{आपका आईटी विभाग/प्रकोष्ठ वर्तमान में कौन से आईटी उपकरण, प्लेटफ़ॉर्म या सॉफ्टवेयर संचालित करता है?}}
\ans\ans\ans

\q{B2}{What mechanisms does the IT department use to support other employees in adopting and using digital tools?\\
\hindi{अन्य कर्मचारियों को डिजिटल उपकरण अपनाने में आपका विभाग क्या सहायता देता है?}}
\choices{
  \item Regular training workshops \hindi{(नियमित प्रशिक्षण)}
  \item On-call helpdesk support \hindi{(हेल्पडेस्क)}
  \item On-site troubleshooting visits \hindi{(साइट पर समस्या समाधान)}
  \item User manuals / documentation \hindi{(उपयोगकर्ता पुस्तिका)}
  \item No formal support mechanism exists \hindi{(कोई औपचारिक सहायता नहीं)}
  \item Other: \rule{5cm}{0.3pt}
}

\q{B3}{Does your department have any internally developed tools, dashboards, or platforms custom-built for your specific workflows?\\
\hindi{क्या आपके विभाग के पास आंतरिक रूप से विकसित कोई उपकरण, डैशबोर्ड या प्लेटफ़ॉर्म है?}}
\choices{
  \item Yes -- describe: \rule{8cm}{0.3pt}
  \item No, we only use centrally provided tools \hindi{(केवल केन्द्रीय उपकरणों का उपयोग)}
}

\q{B4}{In your assessment, what percentage of employees in this department actively use the IT tools provided to them?\\
\hindi{आपके अनुमान अनुसार इस विभाग में कितने प्रतिशत कर्मचारी दिए गए आईटी उपकरणों का सक्रिय उपयोग करते हैं?}}
\choices{
  \item Less than 25\% \quad
  \item 25--50\% \quad
  \item 50--75\% \quad
  \item More than 75\%
}

\q{B5}{What are the biggest barriers to IT adoption that you observe among your department's employees?\\
\hindi{आपके विभाग में आईटी अपनाने में सबसे बड़ी बाधाएं क्या हैं?}}
\ans\ans

\q{B6}{What additional resources or policy changes would most improve IT adoption in your department?\\
\hindi{आईटी अपनाने में सुधार के लिए किन संसाधनों या नीतिगत परिवर्तनों की सबसे अधिक आवश्यकता है?}}
\ans\ans

\end{tcolorbox}

\vspace{0.4cm}

%===== GAP IDENTIFICATION BOX (for surveyor use only) =====
\begin{tcolorbox}[colback=white,colframe=govgold,arc=2mm,boxrule=0.6pt,
  title={\textbf{Surveyor's Gap Classification (Fill after interview -- not shown to respondent)}},fonttitle=\small\bfseries]
\small
Based on this respondent's answers, the primary gap observed is:
\begin{multicols}{2}
\choices{
  \item Skill Gap (trained but not competent)
  \item Awareness Gap (unaware tool exists)
  \item Availability Gap (no hardware/internet)
  \item Motivation Gap (prefers manual method)
  \item Process Gap (tool doesn't fit workflow)
  \item No Gap (fully adopting IT tools)
}
\end{multicols}
\noindent\textbf{Surveyor notes:}\\
\ans\ans
\end{tcolorbox}

\vspace{0.2cm}
\begin{center}
\small\textit{Thank you for your time. \hindi{आपके समय के लिए धन्यवाद।}}
\end{center}

\end{document}
